
\textbf{Proposition 5.2.} \uline{On a sur} \(P\) \uline{la suite exacte} :
\[
0\longrightarrow\underline{\Omega}^1_{P/Y}\longrightarrow f^*(\underline{E})(-1)\longrightarrow\underline{O}_P\longrightarrow0,
\]
\uline{i}.\uline{e}. \uline{on a un isomorphisme}
\[
\underline{\Omega}^1_{P/Y}\simeq \underline{F}(-1),\quad\text{où}\quad \underline{F}=\Ker\bigl(f^*(\underline{E})\longrightarrow\underline{O}_P(1)\bigr).
\]

On considère le fibré vectoriel sur \(Y\) défini par \(\underline{E}\), et on note \(Z\) le complémentaire de la section nulle de ce fibré. On a le diagramme commutatif :
\[
\begin{tikzcd}[column sep=large,row sep=large]
P \arrow[dr,"f"'] && Z \arrow[ll,"h"'] \arrow[dl,"g"]\\
& Y. &
\end{tikzcd}
\]
Le morphisme \(h\) identifie \(Z\) au \(P\)-schéma affine défini par le faisceau d'algèbres \(\underline{B}=\bigoplus_{n\in\mathbb Z}\underline{O}_P(n)\) ; \(\underline{B}\) est une \(\underline{O}_P\)-Algèbre graduée.

Comme \(f,g,h\) sont des morphismes lisses, on peut écrire la suite exacte des faisceaux de différentielles relatives ; comme \(Z\) est affine sur \(P\), on peut écrire cette suite en termes de \(\underline{B}\)-Modules; on obtient ainsi :
\begin{equation}\tag{5.2.1}
0\longrightarrow
 \underline{\Omega}^1_{P/Y}\otimes_{\underline{O}_P}\underline{B}
\longrightarrow
 \underline{\Omega}^1_{\underline{B}/Y}
\longrightarrow
 \underline{\Omega}^1_{\underline{B}/\underline{O}_P}
\longrightarrow0 .
\end{equation}

Or cette suite exacte est une suite exacte graduée de \(\underline{B}\)-Modules gradués. En effet, elle résulte du diagramme commutatif :
\[
\begin{tikzcd}[column sep=large,row sep=large]
\underline{\underline{S}}(\underline{E}) \arrow[rr] && \displaystyle\bigoplus_{n\in\mathbb Z}\underline{\underline{S}}(\underline{E})(n)\simeq \underline{\underline{S}}(\underline{E})[T,T^{-1}]\\
& \underline{O}_X \arrow[ul] \arrow[ur] &
\end{tikzcd}
\]
où les flèches sont des homomorphismes d'algèbres, compatibles avec la graduation de \(\underline{\underline{S}}(\underline{E})[T,T^{-1}]\) (graduation définie par le degré par rapport à l'indéterminée \(T\)), les deux autres algèbres étant supposées réduites
